\documentclass[11pt]{article}
\newcommand{\SheetCode}{Computation C1}
\usepackage[margin=25mm]{geometry}
\usepackage[T1]{fontenc}
\usepackage{lmodern}
\DeclareUnicodeCharacter{2605}{$\star$}
\usepackage{microtype}
\usepackage{amsmath,amssymb,mathtools,bm}
\usepackage{enumitem}
\usepackage{xcolor}
\usepackage{fancyhdr}
\usepackage{hyperref}
\usepackage[most]{tcolorbox}
\definecolor{agiBlue}{HTML}{173B57}
\definecolor{agiGold}{HTML}{B7791F}
\definecolor{agiLight}{HTML}{EFF6FA}
\definecolor{agiGreen}{HTML}{184D3B}
\hypersetup{colorlinks=true,linkcolor=agiBlue,urlcolor=agiBlue}
\setlength{\parindent}{0pt}
\setlength{\parskip}{0.55em}
\setlength{\headheight}{14pt}
\setlist{nosep,leftmargin=*}
\pagestyle{fancy}
\fancyhf{}
\fancyhead[L]{\small Part of the \href{https://github.com/mmikhasenko/GIModel.jl}{Agentic Gottried--Isgur} project}
\fancyhead[R]{\SheetCode}
\fancyfoot[C]{\thepage}
\newcommand{\dd}{\,\mathrm d}
\newcommand{\ket}[1]{\lvert #1\rangle}
\newcommand{\bra}[1]{\langle #1\rvert}
\newcommand{\braket}[2]{\langle #1\mid #2\rangle}
\newcommand{\expect}[1]{\left\langle #1\right\rangle}
\newcommand{\op}[1]{\widehat{#1}}
\newcommand{\GeV}{\,\mathrm{GeV}}
\newcommand{\R}{\mathbb{R}}
\newtcolorbox{goals}{colback=agiLight,colframe=agiBlue,title=Learning outcomes}
\newtcolorbox{reading}{colback=white,colframe=agiGold,title=Book reference,
  after skip=5mm}
\newtcolorbox{paperbridge}{colback=agiGreen!7,colframe=agiGreen,
  colbacktitle=agiGreen,coltitle=white,title=Connection to the GI paper}
\newtcolorbox{problem}[1]{colback=white,colframe=agiBlue,title={#1},
  before skip=3.5mm,after skip=1mm}
\newtcolorbox{solution}{colback=black!2,colframe=black!45,title=Solution}
\newtcolorbox{quiz}{colback=white,colframe=agiGold,title=Post-solution concept check}
\newtcolorbox{checkpoint}{colback=white,colframe=agiGold,title=Interpretation checkpoint}
\newcommand{\sheettitle}[3]{%
  \begin{center}
  {\Large\bfseries #1}\par
  \vspace{0.2em}{\color{agiBlue}#2}\par
  \vspace{0.4em}\small #3
  \end{center}\vspace{0.5em}}
\newcommand{\difficulty}[1]{\textbf{Difficulty:} #1}
\newcommand{\timebox}[1]{\textbf{Suggested time:} #1}
\newcommand{\answer}{\par\smallskip\color{black!50}\textbf{Answer:} }

\begin{document}
\sheettitle{Computational discovery map}{Questions before implementation}{Prerequisite: Theory 01--08; Theory 09 recommended. \difficulty{★/★★★} \quad \timebox{2 hours reading + exploration}}

\begin{goals}
Turn the theoretical dependency graph into numerical questions; recognize the
minimum tests for a trustworthy implementation; prepare a later from-scratch
Julia course without coupling the physics to one representation.
\end{goals}

\begin{checkpoint}
This is intentionally a question-led sheet, not a coding recipe. First answer
each question in plain language. For a proposed test, give its input,
expected result, and failure criterion; for a procedure, give brief pseudocode.
Numerical choices below are toy inputs, not GI parameter fits.
Then locate the named repository component.
In the later hands-on course, implement the smallest version from scratch before
reading the production implementation.
\end{checkpoint}

\section*{Level 1 — Numbers and units}
\begin{enumerate}
\item ★ In natural units, what are the dimensions of $r$, $b$, $\sigma$, and a
radial reduced wave $u(r)$? Use $\hbar=c=1$, $V(r)=br$,
$\rho_\sigma\propto e^{-\sigma^2r^2}$, and $\int|u|^2\dd r=1$.
Which invalid expression could a unit-aware type
system reject?
\item ★ How many decimal digits are lost when subtracting two nearly equal
energies to report a fine-structure splitting? Use $E_1=3.100000\,\mathrm{GeV}$
and $E_2=3.100010\,\mathrm{GeV}$ as an example, and compare
$\max(|E_1|,|E_2|)/|E_2-E_1|$. Assume comparable absolute input errors.
When should the working precision exceed the output precision?
\item ★ Why should the normalization use quadrature weights instead of the raw
Euclidean norm of a sample vector? Let $u_i=u(r_i)$, with positive weights
$w_i$ approximating $\int f(r)\dd r$ by $\sum_iw_if(r_i)$.
Write the discrete norm and its normalization factor. Inspect the contracts in
\texttt{src/model\_objects.jl}.
\item ★ What absolute and relative tolerances are sensible for a quantity that
can legitimately cross zero? Use the comparison
$|x-y|\le a_{\rm tol}+r_{\rm tol}\max(|x|,|y|)$.
For a splitting with target accuracy $10^{-5}\,\mathrm{GeV}$, propose
$a_{\rm tol}$ with units and a dimensionless $r_{\rm tol}$; explain their roles
near zero and near $1\,\mathrm{GeV}$. Why is one global tolerance dangerous?
\end{enumerate}

\section*{Level 2 — Two representations of one radial problem}
\begin{enumerate}[resume]
\item ★ Derive the three-point discretization of $-u''$ on $r_i=ih$. Where do
the boundary conditions enter? Take $i=1,\ldots,N$,
$u_0=u_{N+1}=0$, and $r_{\max}=(N+1)h$; expand $u(r_i\pm h)$ through
fourth order and state the truncation error. Compare with \texttt{src/radial\_grid.jl}.
\item ★★ If the discrete $p^2$ matrix is diagonalized as $Q\Lambda Q^T$, how do
you compute $\sqrt{p^2+m^2}$ without taking elementwise square roots?
Assume a real symmetric nonnegative $p^2$ matrix in an orthonormal
representation, $Q^TQ=I$, and $m>0$. Write the matrix expression.
\item ★ Write the first three radial HO basis functions for fixed $L$. Which
normalization and phase tests would you require before using them in mixing?
Use Sheet 3, Problem 5, with $n=0,1,2$ and
$L_2^\alpha(x)=[x^2-2(\alpha+2)x+(\alpha+1)(\alpha+2)]/2$.
Inspect \texttt{src/harmonic\_oscillator\_basis.jl}.
\item ★★ Design independent refinement paths for a grid $(h,r_{\max})$ and an
HO basis $(N,\beta)$. What false plateau can occur if only one control changes?
\item ★★★ What quantities can be compared across HO and FD without converting
one stored wave into the other's array type? Use the interface described in
\texttt{docs/code\_architecture.md}.
\end{enumerate}

\section*{Level 3 — Operators, not formulas}
\begin{enumerate}[resume]
\item ★ Implement on paper a matrix function $f(p^2)=Qf(\Lambda)Q^T$. What
properties of $f$ and $p^2$ make the result real symmetric?
\item ★★ Given matrices $B$ and $V$, how would you test $BVB$ for Hermiticity,
scale-invariant expectations, and convergence separately?
Assume $B,V$ are Hermitian in an orthonormal basis. For a nonzero vector $v$,
use $\langle O\rangle_v=v^\dagger Ov/(v^\dagger v)$; compare $v$ and $3v$.
For convergence, specify a sequence of increasing basis sizes.
\item ★★ Why should analytic derivatives of $\alpha_s(r)$ be tested against,
but not replaced by, finite differences? Inspect
\texttt{src/running\_coupling.jl} and \texttt{src/spin\_fine\_structure.jl}.
\item ★★★ Design a test that would detect confusion between $R(r)$ and $u(r)$
even if all eigenvectors were normalized to one.
\end{enumerate}

\section*{Level 4 — Eigensystems and state identity}
\begin{enumerate}[resume]
\item ★ How do you sort eigenpairs while preserving the correspondence between
eigenvalues and vectors? What additional label distinguishes two equal-$J$
states before mixing?
\item ★★ Near an avoided crossing, why can ``state number 1'' cease to track the
same physical component? Consider the toy matrix
$H(t)=\begin{pmatrix}t&v\\v&-t\end{pmatrix}$ with fixed $v>0$,
where $t,v$ have energy units. Compare ordered energies and basis-component
weights as $t$ crosses zero. Propose a rule using overlaps between normalized
eigenvectors at neighboring $t$, and distinguish following an eigenstate
continuously from following its dominant basis component.
\item ★★ What must a convergence certificate record so another person can audit
why a solve stopped? Inspect \texttt{src/solver\_options.jl}.
\item ★★★ Why should a fixed-sector Hamiltonian be solved once, with all its
reported contributions evaluated in that eigenstate, rather than solving four
slightly different Hamiltonians?
\end{enumerate}

\section*{Level 5 — Mixing and observables}
\begin{enumerate}[resume]
\item ★ Implement the analytic checks $\operatorname{tr}M=\sum\lambda_i$ and
$\det M=\prod\lambda_i$ for every $2\times2$ mixing block. What do they not test?
\item ★★ How will you make eigenvector phases deterministic for display while
keeping tests phase-invariant? For a normalized complex vector $c$, propose
a phase choice based on a nonzero component and a rule for ties or nearly zero
components. Which overlap or projector would you compare in a test? Inspect \texttt{src/state\_mixing.jl}.
\item ★★★ Design a data structure for a physical state mixed first in
spectroscopic space and then in flavor space. How do you retain each component
wave without inventing a representative wave?
\item ★★ Which observables stress the tail, the bulk, and the origin of a wave?
Consider $\expect{r^2}$, the energy, and $|R(0)|^2$.
Assign each to the region it emphasizes and choose one convergence test for each.
\end{enumerate}

\section*{Level 6 — Reproducibility as an algorithm}
\begin{enumerate}[resume]
\item ★ Why must digitized paper values live outside the pure model package?
Inspect the boundary between \texttt{src/} and \texttt{GIPaper/}.
\item ★★ Design a residual table that cannot hide one bad flavor sector behind
a good global average. Specify the columns for state labels, flavor,
calculated and reference masses, signed residuals, and source; state how missing
reference values are recorded. A blank reference must not count as zero residual.
\item ★★ Which inputs require provenance, units, confidence, and an explicit
``unknown'' state rather than a guessed fallback?
\item ★★★ Draft the smallest continuous-integration gate that tests data,
physics, both solvers, and one end-to-end reference comparison without making
every commit prohibitively slow.
\end{enumerate}

\section*{Proposed later hands-on sequence}
The answers naturally become six small implementations: (1) normalized radial
grid and Coulomb solver; (2) HO Rayleigh--Ritz solver; (3) matrix functions and
the semirelativistic kinetic operator; (4) smearing plus $BVB$ operators; (5)
fixed sectors and $2\times2$ mixing; (6) an observable and a provenance-aware
comparison report. Each unit should begin with an analytic toy case and end with
an HO--FD or invariant cross-check.

\begin{quiz}
Exit question: if a program reproduces the target masses but fails normalization,
Hermiticity, or convergence tests, has it reproduced GI? \answer no. It has fit
numbers without establishing the computation that produced them.
\end{quiz}
\end{document}
